% Generated draft: requires mathematical and source review.
\documentclass[11pt]{article}
\usepackage[margin=25mm]{geometry}
\usepackage{fontspec}
\setmainfont{Latin Modern Roman}
\usepackage{amsmath,amssymb,mathtools,unicode-math}
\setmathfont{Latin Modern Math}
\usepackage{enumitem,needspace,xurl,hyperref}
\setlength{\parindent}{0pt}
\setlength{\parskip}{5pt}
\setlength{\emergencystretch}{3em}
\newcommand{\N}{\mathbb{N}}
\newcommand{\Z}{\mathbb{Z}}
\newcommand{\Q}{\mathbb{Q}}
\newcommand{\R}{\mathbb{R}}
\newcommand{\C}{\mathbb{C}}
\newcommand{\F}{\mathbb{F}}
\newcommand{\A}{\mathbb{A}}
\newcommand{\PP}{\mathbb{P}}
\newcommand{\E}{\mathbb{E}}
\newcommand{\eps}{\varepsilon}
\newcommand{\dd}{\,\mathrm{d}}
\begin{document}
\section*{UnsolvedMath ID: 2307022}
\textbf{Source problem code:} catalogue entry
\subsection{Definitions and mathematical statement}
A Jordan curve in \(\R^3\) is the image of an injective continuous map from \(S^1\). For disjoint Jordan curves \(A,B\), define
\[
\operatorname{dist}(A,B)=\inf\{\|a-b\|:a\in A,\ b\in B\}.
\]
Let \(\operatorname{len}(A)\) be its arc length, interpreted as \(\infty\) when it is nonrectifiable. The assertion attributed to Gehring is
\[
\forall A,B\text{ disjoint linked Jordan curves in }\R^3,
\qquad\operatorname{dist}(A,B)=1\ \Longrightarrow\
\operatorname{len}(A)\ge2\pi.
\]
Only the length of \(A\) is bounded; no length or curvature constraint on \(B\) is supplied. The source credits Gehring with the corresponding inequality using some positive absolute constant in place of \(2\pi\).

The meaning of ``linked'' is not defined in the excerpt. Nonzero linking number and nonsplitting (no separating embedded sphere) are different notions; the latter also includes some links of linking number zero. No choice between these is silently imposed.

The structured registry status is partially solved with classification PARTIAL-PROGRESS, whereas its research summary asks to mark SOLVED-IN-LITERATURE with moderate confidence. That summary asserts the sharp \(2\pi\) bound but gives no identifiable theorem or citation, and its ``up to the constant'' wording does not state a precise exact-constant result. Accordingly the exact \(2\pi\) target is retained without promoting this vague resolution claim to a verified theorem.
\subsection{Short English statement}
Two linked disjoint closed curves at mutual distance one should each have length at least \(2\pi\). The linking convention and the supplied inconsistent, uncited resolution assessment require review.
\subsection{Sources}
\begin{itemize}
\item Local UnsolvedMath registry, permanent ID \(2307022\), exact source code \texttt{catalogue entry}, ``Research Problems in Function Theory --- Problem 7.22.''
\item Hayman, \emph{Research Problems in Function Theory} (2018), Problem 7.22, attributed to F. W. Gehring; \url{https://arxiv.org/abs/1809.07200}. Source-TeX extraction; recorded access 2026-07-29. The Gehring positive-constant result has no complete citation here.
\item Supplied research summary refers vaguely to work on the length of linked curves, without an authoritative paper, author, theorem number, or exact-constant statement supporting its moderate-confidence solved label.
\end{itemize}
\end{document}
